We represent simultaneous encounters through a local active conflict graph. 
For each sensing edge $\{i,j\}\in\mathcal E_N(t)$, define the nominal pairwise CBF margin
\begin{equation}
    \sigma_{ij}^{\mathrm{nom}}=
    L_{f_{ij}}h_{ij}+
    L_{g_i}h_{ij}u_i^{\mathrm{nom}}+
    L_{g_j}h_{ij}u_j^{\mathrm{nom}}+
    \alpha_c(h_{ij}).
    \label{eq:nominal_pairwise_margin}
\end{equation}
An edge becomes active when $\sigma_{ij}^{\mathrm{nom}}\leq\eta_{\mathrm{on}}$ and remains active until $\sigma_{ij}^{\mathrm{nom}}\geq\eta_{\mathrm{off}}$, where $0\leq\eta_{\mathrm{on}}<\eta_{\mathrm{off}}$.
The resulting active graph is denoted by $\mathcal G_a(t)=(\mathcal A,\mathcal E_a(t))$.

When robot $i$ first becomes incident to an active edge, it freezes the node-wise precedence key
\begin{equation}
    \kappa_i=(\bar\mu_i,-i),
    \label{eq:precedence_key}
\end{equation}
where $\bar\mu_i$ is its current energy margin and the robot identifier provides deterministic tie-breaking.
The same key is used for all active neighbors until robot $i$ has no remaining active edge.
Each edge is directed from the robot with the larger key toward the robot with the smaller key, where the edge direction denotes yielding.
Because all edges follow the same strict total ordering, the directed active graph is acyclic.

Precedence also selects the reference input of the safety filter:
\begin{equation}
    u_i^{\mathrm{ref}}
    =
    \begin{cases}
        u_i^{\mathrm{nom}},
        & \text{if robot $i$ has no outgoing edge},\\
        u_i^{\mathrm{yield}},
        & \text{otherwise}.
    \end{cases}
    \label{eq:precedence_reference}
\end{equation}
Here, $u_i^{\mathrm{yield}}$ directs the robot toward a locally available holding region outside the conflict zone.

For an active pair, let $y_{ij}$ and $r_{ij}$ denote the yielding and right-of-way robots, respectively, and define
\begin{equation}
    a_{ij}^{s}(u_s)=
    L_{g_s}h_{ij}u_s+
    \frac{1}{2}
    \left(
        L_{f_{ij}}h_{ij}
        +
        \alpha_c(h_{ij})
    \right),
    \; s\in\{i,j\}.
    \label{eq:local_pairwise_contribution}
\end{equation}
The responsibility transfer is
\begin{equation}
\begin{aligned}
    \rho_{ij}
    =
    \min\left\{
        \left[-a_{ij}^{r_{ij}}
        (u_{r_{ij}}^{\mathrm{ref}})\right]_+,\,
        \left[\max_{u\in\mathcal U}
        a_{ij}^{y_{ij}}(u)\right]_+
    \right\},
\end{aligned}
\label{eq:responsibility_transfer}
\end{equation}
and the complementary allocations are
\begin{equation}
    \gamma_{y_{ij}r_{ij}}=-\rho_{ij},
    \qquad
    \gamma_{r_{ij}y_{ij}}=\rho_{ij}.
    \label{eq:complementary_responsibility}
\end{equation}
For inactive pairs, $\gamma_{ij}=0$.

Robot $i$ then solves
\begin{equation}
\begin{aligned}
    u_i^\star
    =
    \arg&\min_{u_i\in\mathcal U}
    \quad
    \|u_i-u_i^{\mathrm{ref}}\|^2\\
    \mathrm{s.t.}\quad&
    L_{g_i}h_{ij}u_i
    +
    \frac{1}{2}
    \left(
        L_{f_{ij}}h_{ij}
        +
        \alpha_c(h_{ij})
    \right)
    +
    \gamma_{ij}
    \geq0,\\
    &
    L_{f_i}h_{i,o}^{\mathrm{obs}}
    +
    L_{g_i}h_{i,o}^{\mathrm{obs}}u_i
    \geq
    -\alpha_o(h_{i,o}^{\mathrm{obs}}),
    \quad \forall o\in\mathcal O_i,\\
    &
    L_{F_i}h_{e,i}
    +
    L_{G_i}h_{e,i}u_i
    \geq
    -\alpha_e(h_{e,i}).
    \label{eq:distributed_safety_filter}
\end{aligned}
\end{equation}

\begin{proposition}[Local precedence consistency and safety]
Suppose that neighboring robots maintain consistent node-wise keys, the local QPs remain feasible, and discrete certificate updates preserve $h_{e,i}\geq0$.
Then the directed active graph is acyclic, the filtered inputs satisfy the centralized pairwise collision-avoidance condition, and the proxy energy-sufficiency sets are preserved between certificate updates.
When the calibrated station-cost event holds, the station-arrival voltage is no lower than $V_{\min}$.
\end{proposition}

Acyclicity follows from the strict ordering of the node-wise keys.
Moreover, $\gamma_{ij}+\gamma_{ji}=0$, so summing the paired local constraints recovers the centralized pairwise CBF condition.
The acyclic ordering prevents cyclic precedence among locally observed robots but does not guarantee global deadlock freedom.
Progress additionally requires a feasible local holding region and feasibility of the combined CBF constraints.
